\chapter{量子力学の完成とAI --- 確率の世界と「次元の呪い」の打破}

\section*{本章の学習目標}
\begin{itemize}
\item 物質も「波」であるというド・ブロイの物質波の概念を理解する。
\item 「物理量」が「オペレーター（演算子）」に変わる量子力学の数学的構造を知る。
\item シュレディンガー方程式における「時間の1階微分と位置の2階微分」、および「虚数\(i\)」が不可欠な理由を知る。
\item オペレーターから飛び出す「固有値」が現実の観測データであることを理解する。
\item ハイゼンベルクの行列力学と不確定性原理の真の理由（非可換性）を学ぶ。
\item 波束の収縮、シュレディンガーの猫、量子もつれなど、確率解釈がもたらした決定論の終焉を理解する。
\item 量子力学に立ちはだかる「次元の呪い（スパコンの限界）」と、それを打破するAIの最前線を知る。
\end{itemize}

\section{物質も「波」である：ド・ブロイの逆転の発想}

第8章の前期量子論で見たように、アインシュタインは「光は波であると同時に、粒子（フォトン）でもある」という光量子仮説で世界を驚かせました。

\subsection{光の粒子性を決定づけた「コンプトン効果」}

このアインシュタインの予言を、疑いようのない形で実証したのが1923年のコンプトン効果の発見です。

アメリカの物理学者アーサー・コンプトンは、X線（波長が短くエネルギーの高い光）を物質中の電子にぶつけ、散乱されたX線の波長を測定する精密な実験を行いました。

\subsection{発想の背景と波の理論での苦闘：}

当時、X線やガンマ線を物質に当てて散乱させると、なぜか跳ね返ってきた線が「軟らかくなる（物を突き抜ける力が弱くなる＝波長が長くなる）」という不可解な現象が、複数の研究者によって報告されていました。

もしX線がマクスウェルの言うような純粋な「波」であるならば、X線の波が電子に当たると、電子はその波の揺れ（振動数）に合わせて同じペースで揺さぶられ、受け取ったのと同じ振動数（＝同じ波長）のX線を四方八方に再放射するはずです。つまり、古典物理学では\textbf{「散乱される前後でX線の波長は絶対に変わらない」}と予測されていました（トムソン散乱）。

コンプトン自身も、この波長が伸びる謎を解明しようと数年間がかりで実験を続けていました。彼は当初、光は「波」であるという常識を固く信じており、なんとか古典的な波の理論（ドップラー効果など）を使ってこの現象を説明しようと必死に数式をこねくり回しましたが、どうやっても実験データと辻褄が合いませんでした。

\subsection{「粒子」へのパラダイムシフト：}

行き詰まったコンプトンは、ついに視点を180度変える決断をします。当時まだ物理学界から異端視され、ほとんど誰も信じていなかったアインシュタインの\textbf{「光量子仮説（光は粒である）」}を、あえて使ってみることにしたのです。

コンプトンは、光を「エネルギー\(E=h\nu\)、運動量\(p=h/\lambda\)を持つビリヤードの球」とみなして、静止している電子（もう一つの球）に激突させる計算を行いました。衝突によって光の粒（フォトン）は電子を弾き飛ばし（エネルギーと勢いを与え）、自らはエネルギーを失って斜めに跳ね返ります。

フォトンがエネルギーを失うということは、\(E=h\nu\)より振動数\(\nu\)が小さくなる、すなわち「波長\(\lambda\)が長くなる」ことを意味します。

コンプトンは、高校物理で習う「エネルギー保存の法則」と「運動量保存の法則」を、このフォトンと電子の衝突にそのまま適用して計算し、波長のズレ（\(\Delta \lambda\)）を表す次の方程式を導き出しました。

\[\Delta \lambda = \lambda' - \lambda = \frac{h}{m_e c} (1 - \cos\theta)\]

（\(\lambda\)は衝突前の波長、\(\lambda'\)は衝突後の波長、\(m_e\)は電子の質量、\(\theta\)はX線が跳ね返った角度）

この「ビリヤードの衝突計算」から導かれた波長のズレは、彼自身の精密な実験データを高精度に説明しました。波であるはずの光が、空間を飛ぶ「粒」のように運動量（勢い）を持ち、電子と衝突することが強く示されたのです。（※オランダのピーター・デバイもほぼ同時期に独立して同じ理論に到達しました）。この決定的な証拠により、物理学界はついにアインシュタインの光量子仮説を受け入れることになります。

\subsection{逆転の発想：物質の「波」}

光が粒子としての性質を確固たるものにしたまさにその翌年（1924年）、フランスの若き貴族出身の大学院生、ルイ・ド・ブロイは、さらに常識外れな逆転の発想を思いつきます。「自然界は対称性を好むはずだ。もし波だと思われていた光が粒子の性質を持つのなら、逆に\textbf{『粒子』だと思われている電子や原子も、実は『波』の性質を持っているのではないか？}」

\subsection{光の数式から物質の数式への飛躍}

ド・ブロイは、この突飛なアイデアを数学的に裏付けるため、アインシュタインの2つの偉大な発見（特殊相対性理論と光量子仮説）を組み合わせました。

第5章で見たように、質量ゼロの光は\(E = pc\)というエネルギーと運動量（\(p\)）の関係を持ちます。これを変形すると\(p = E/c\)となります。

そして第8章で見たように、光の粒のエネルギーは\(E = h\nu\)（プランク定数\(\times\)振動数）です。

さらに、波の基本的な性質として、光の速さは「波長\(\times\)振動数（\(c = \lambda\nu\)）」で表されます。

これらをパズルのように組み合わせると、光の運動量\(p\)は次のように計算できます。

\[p = \frac{E}{c} = \frac{h\nu}{\lambda\nu} = \frac{h}{\lambda}\]

これを波長（\(\lambda\)）について解き直すと、\(\lambda = \frac{h}{p}\)となります。

つまり「光の波長は、プランク定数を運動量で割ったものである」という関係が導き出されたのです。

ド・ブロイの天才的な飛躍は、\textbf{「光について成り立ったこの美しい数式が、質量を持つ普通の物質（電子など）にもそのまま当てはまるはずだ」と考えたことでした。 彼はあらゆる物質が持つこの波を「物質波（ド・ブロイ波）」}と名付けました。質量\(m\)を持ち、速度\(v\)で走っている物質（つまり運動量\(p=mv\)を持つ物質）は、次のような波長（\(\lambda\)）の波として振る舞うと宣言したのです。

\[\lambda = \frac{h}{p} = \frac{h}{mv}\]

例えば、野球のボールが飛んでいるときも、この数式によればボールは「波」として揺れながら飛んでいることになります。しかし、ボールは質量\(m\)が大きすぎるため、分母が巨大になって波長\(\lambda\)があまりにも短くなりすぎ、波としての性質は完全に隠れてしまいます。ところが、質量が極めて小さい「電子」などでは、波長が十分に長くなり、波の性質がはっきりと現れるのです。

ド・ブロイのこの突飛な論文は審査員たちを困惑させましたが、論文を送られたアインシュタインが「これは天才のひらめきだ」と絶賛したことで、物理学の歴史は大きく動き出します。

\subsection{物質が波になる瞬間：親子の皮肉な歴史と電子線回折}

ド・ブロイの「物質波」の予言は、そのわずか3年後（1927年）、2つの独立した電子線の回折実験によって見事に実証されました。

一つは、アメリカの物理学者クリントン・デヴィソンとレスター・ガーマーによる実験です。彼らがニッケルの結晶に向けて電子の「粒」のビームを撃ち込んだところ、跳ね返った電子はランダムに散らばるのではなく、特定の方向にだけ集中して飛んでくることが分かりました。

もう一つは、イギリスの物理学者ジョージ・パジェット・トムソンによる、金属の薄膜に電子線を透過させる実験です。彼もまた、X線（光の波）を当てたときと同じような美しい同心円状の干渉縞（波特有の模様）を、電子のビームを使って写真乾板上に撮影することに成功しました。

これらはまさしく、波と波が重なり合って強め合う「干渉（回折）」という波特有の現象でした。実験から逆算された電子の波長は、ド・ブロイの数式が示す値とよく一致したのです。（この功績により、デヴィソンとG.P.トムソンは1937年にノーベル物理学賞を共同受賞します）。

ここで、物理学の歴史における極めて面白く、ドラマチックな皮肉を紹介しましょう。

この実験を成功させたG.P.トムソンは、1897年に「電子が『粒子』であることを世界で初めて発見した」J.J.トムソンの実の息子だったのです。つまり、\textbf{「父親が電子を『粒子』だと証明してノーベル賞を獲り、その息子が同じ電子を『波』だと証明してノーベル賞を獲った」}という歴史的快挙でした。

コンプトン効果によって「波（光）が粒である」ことが証明され、トムソン親子の歴史をまたいだ実験によって「粒（電子）が波である」ことが証明される。この見事な対称性により、「波と粒子の二重性」という量子力学の最も不思議な大前提が完全に確立されました。

\section{シュレディンガー方程式と「虚数」の神秘}

ド・ブロイの「電子は波である」というアイデアに強く感銘を受けたのが、オーストリアの物理学者エルヴィン・シュレディンガーでした。

「もし電子が本当に波ならば、光や音の波と同じように、その動きを計算できる『波動方程式』が必ず存在するはずだ」。彼は1926年、アルプスでの休暇中に猛烈な計算に没頭しました。

\subsection{物理量が「数字」から「オペレーター（装置）」へ変わる}

第1章で見たように、マクロな世界ではニュートンの「運動方程式（\(F=ma\)）」を使えば、ボールの未来の動きを計算できました。しかし、電子は波の性質を持つため、ニュートン力学の計算式がそのままでは全く使えません。

ここでシュレディンガーは、物理学における極めて重大なパラダイムシフトとなる数学的飛躍を行います。それは、\textbf{「物理量を、単なる『数字』から、波を加工する『オペレーター（演算子：数学的な装置）』に置き換える」}という手法でした。

古典力学では、運動量\(p\)やエネルギー\(E\)は「5」や「10」といったただの数字です。しかし量子力学では、物質は波（波動関数\(\psi\)）として広がっているため、ある特定の「数字」を持ちません。その代わり、\textbf{「波の形を調べて（微分して）、そこに隠されている運動量やエネルギーの値を引きずり出すための装置（オペレーター）」}を波に作用させるのです。

\subsection{平面波からの「逆算」：オペレーターはいかにして誕生したか}

シュレディンガーがこの「微分オペレーター」という魔法のようなアイデアに行き着いたのは、\textbf{「何もない空間を飛んでいる電子（自由粒子）は、最もシンプルで綺麗な『平面波』として振る舞うはずだ」}という大前提からスタートしたためです。

彼はまず、ド・ブロイの波長（\(p = h/\lambda\)）とアインシュタインの光量子仮説（\(E = h\nu\)）に従って飛んでいる電子の波を、オイラーの公式を使って次のような美しい複素数の平面波として書き下しました。（※\(\hbar\)はプランク定数\(h\)を\(2\pi\)で割ったものです）

\[\psi = A e^{i(px - Et)/\hbar}\]

シュレディンガーはこの波の式（\(\psi\)）を眺めていて、決定的な事実に気づきます。

「この波を空間\(x\)で微分すれば、肩に乗っている『運動量\(p\)』が前に飛び出してくるじゃないか」

\[\frac{\partial \psi}{\partial x} = \frac{ip}{\hbar} \psi\]

これを\(p\)について解き直す（両辺に\(-i\hbar\)を掛ける）と、次のようになります。

\[p \psi = \left( -i\hbar \frac{\partial}{\partial x} \right) \psi\]

ここから、波から運動量を取り出すための装置として \textbf{「運動量オペレーター（\(\hat{p} = -i\hbar \frac{\partial}{\partial x}\)）」}を思いつきました。

同様に、\textbf{「時間\(t\)で微分すれば、肩に乗っている『エネルギー\(E\)』が飛び出してくる」}ことにも気づきます。

\[\frac{\partial \psi}{\partial t} = -\frac{iE}{\hbar} \psi\]

これを\(E\)について解き直すと、次のようになります。

\[E \psi = \left( i\hbar \frac{\partial}{\partial t} \right) \psi\]

ここから、波からエネルギーを取り出す装置として \textbf{「エネルギーオペレーター（\(\hat{E} = i\hbar \frac{\partial}{\partial t}\)）」}を導き出したのです。

シュレディンガーは、「この最も単純な平面波から逆算して作った『微分ルール（オペレーター）』は、水素原子の周りのような複雑な空間で波が歪んで定在波になっていても、そのまま汎用的に使える絶対的なルールに違いない」という、極めて数学的で大胆な飛躍を行いました。

\subsection{驚くべき発想：「時間の1次微分」と「位置の2次微分」の釣り合い}

次にシュレディンガーは、古典力学における最も基本的な「エネルギー保存則」の数式を、波の世界に翻訳しようと試みました。

ゆっくり動く物体の総エネルギー（\(E\)）は、「運動エネルギー（\(\frac{1}{2}mv^2\)、あるいは運動量\(p\)を使って\(\frac{p^2}{2m}\)）」と「位置エネルギー（\(V\)）」の足し算で表されます。

\[E = \frac{p^2}{2m} + V(x)\]

このお馴染みの数式に、先ほどの「オペレーター（微分装置）」をそっくりそのまま当てはめ、電子の波（\(\psi\)）に作用させてみたのです。

左辺のエネルギー\(E\)には「時間の1回微分（\(\frac{\partial}{\partial t}\)）」を当てはめます。

一方、右辺の運動量\(p\)は「2乗（\(p^2\)）」されているため、「空間の微分を2回繰り返す（\(\frac{\partial^2}{\partial x^2}\)）」ことになります。

ここが、シュレディンガーの閃きの凄まじいところでした。

第5章で学んだような、光や音を伝える「古典的な波動方程式（\(\frac{\partial^2 u}{\partial t^2} = v^2 \frac{\partial^2 u}{\partial x^2}\)）」は、時間についても空間についても\textbf{「両方とも2階微分」でした。 しかし、電子の波が力学のエネルギー保存則を満たすためには、「時間の1階微分」と「空間の2階微分」という、アンバランスな微分同士をイコールで釣り合わせなければならない}ことに気づいたのです。

\subsection{なぜ宇宙の法則に「裸の虚数\(i\)」が必要なのか？}

この「時間の1階微分と空間の2階微分を釣り合わせる」という数学的な要請が、宇宙の根本に\textbf{「虚数\(i\)」}を強制的に刻み込むことになります。

波の形をオイラーの公式を使って複素数で表した\(e^{i(kx - \omega t)}\)という波を想像してください。

この波を空間で「2回」微分すると、肩に乗っている\(ik\)が2回前に飛び出してきます。\((ik) \times (ik) = i^2 k^2\)。虚数\(i\)は2乗すると「マイナス1」になるため、虚数が綺麗に消え去り、\textbf{「実数（\(-k^2\)）」}になります。

ところが、この波を時間で「1回」だけ微分すると、肩に乗っている\(-i\omega\)が1回しか飛び出してきません。つまり、結果に\textbf{「虚数\(i\)」がそのまま残ってしまう}のです。

右辺（空間の2回微分）は「実数」なのに、左辺（時間の1回微分）は「虚数」になってしまう。このままでは絶対にイコールで結ぶことはできません。

このアンバランスを解消し、両辺を釣り合わせるための唯一の魔法が、\textbf{「方程式の左辺（時間微分の側）に、あらかじめ『虚数\(i\)』を掛け算しておく」}ことでした。前もって\(i\)を掛けておけば、微分で飛び出してきた\(i\)と合わさって\(i^2 = -1\)となり、両辺が美しい「実数」として釣り合うからです。

\subsection{シュレディンガー方程式の書き下しとハミルトニアン}

こうして、微分オペレーターの前に「虚数\(i\)」を配置するという必然性から、シュレディンガーはミクロの世界を支配する究極の方程式を書き下しました。

これが有名な\textbf{「時間依存のシュレディンガー方程式」}です。

\[i\hbar \frac{\partial \psi}{\partial t} = \left( -\frac{\hbar^2}{2m} \frac{\partial^2}{\partial x^2} + V(x) \right) \psi\]

この方程式の構造を、先ほど定義した「オペレーター」の視点から紐解いてみましょう。

左辺（\(i\hbar \frac{\partial}{\partial t}\)）：これは波を時間で微分してエネルギーを取り出す\textbf{「エネルギーオペレーター（\(\hat{E}\)）」}です。

右辺のカッコ内：空間の2回微分で表される運動エネルギー（\(p^2/2m\)）と、位置エネルギー（\(V(x)\)）を足し合わせたものです。解析力学（第2章）において、システムの全エネルギー（運動エネルギー＋位置エネルギー）を表す概念を\textbf{「ハミルトニアン」と呼びました。量子力学では、この右辺のカッコ全体が全エネルギーを引きずり出す装置、すなわち「ハミルトニアンオペレーター（\(\hat{H}\)）」}となります。

つまり、この方程式を物理的に読むと次のような極めてシンプルな意味になります。

「波を時間で微分して『エネルギー（\(\hat{E}\)）』を取り出した結果は、波を空間で微分して『運動エネルギーと位置エネルギーの合計（\(\hat{H}\)）』を取り出した結果と、常にイコール（同じ）にならなければならない」

これはまさに、古典力学の「エネルギー保存の法則（\(E = K + V\)）」を、波のオペレーター同士の釣り合い（\(\hat{E}\psi = \hat{H}\psi\)）として美しく書き直した数式だったのです。

これは物理学史上、極めて異常な事態でした。これまでの古典物理学の方程式（ニュートンの\(F=ma\)やマクスウェル方程式）はすべて実数だけで書かれており、虚数はあくまで「計算をラクにするための架空のツール」として使われるだけでした。

しかし量子力学においては、\textbf{「虚数\(i\)」は単なる計算ツールではなく、宇宙の根本法則を表す方程式の中に裸のまま組み込まれている「絶対的な構成要素」}だったのです。ミクロの世界の物質は、私たちの目には見えない「虚数を含んだ複素数の波（確率振幅）」として揺らめいていることが確定しました。


\subsection{古典力学では越えられない壁をすり抜ける：量子トンネル効果}

シュレディンガー方程式の威力が最も直感に反して現れる現象の一つが、\textbf{「量子トンネル効果」}です。

古典力学では、粒子が高い壁を越えるためには、その壁を乗り越えるだけのエネルギーが必要です。例えば、ボールを坂の向こう側へ転がしたいなら、ボールには坂の頂上まで登れるだけの運動エネルギーが必要です。エネルギーが足りなければ、ボールは坂の途中で止まり、決して向こう側へ行くことはできません。

ところが、量子力学の世界では話が変わります。電子や原子核のようなミクロな粒子は、単なる小さな球ではなく、波動関数\(\psi\)として空間に広がっています。そのため、粒子のエネルギー\(E\)が、壁の高さに相当する位置エネルギー\(V(x)\)より低い場合でも、波動関数は壁の手前で完全にゼロにはなりません。

古典力学では「侵入禁止」となる領域、すなわち

\[
E < V(x)
\]

となる場所でも、シュレディンガー方程式を解くと、波動関数は壁の中へ少しだけ染み込みます。ただし、その振幅は壁の中で急速に小さくなっていきます。壁が薄ければ、波動関数の一部は反対側までわずかに届きます。そして、反対側で\(|\psi|^2\)がゼロでなければ、そこに粒子が発見される確率もゼロではありません。

これが量子トンネル効果です。

重要なのは、粒子が古典的な意味で「壁に穴を開けて通り抜ける」わけではない、という点です。量子力学では、粒子は最初から波として広がっており、その波が壁の内部で完全には消えず、反対側にもわずかに残るのです。観測したとき、その反対側で粒子が見つかることがある。これが「トンネルを抜けた」ように見える理由です。

この現象は、単なる理論上の奇妙な話ではありません。原子核から\(\alpha\)粒子が飛び出す\(\alpha\)崩壊、太陽内部で起こる核融合、半導体中の電子の移動、さらには原子の表面を観察する走査型トンネル顕微鏡（STM）など、現代科学技術の至るところで量子トンネル効果は実際に働いています。

つまり、シュレディンガー方程式は、ミクロな粒子が「どこにいるか」だけでなく、古典力学では絶対に越えられないはずの境界を、確率の波としてどの程度すり抜けるかまで予測してしまうのです。量子力学における波動関数\(\psi\)は、単なる計算上の記号ではなく、自然が古典的な直感を超えて振る舞うための、極めて強力な数学的表現だったのです。

\subsection{オペレーターから「観測値」を引き出す：固有値と固有関数}

先ほどの方程式の右辺にあるハミルトニアンオペレーター\(\hat{H}\)を使って、もう少し深く見てみましょう。

もし、電子の波（\(\psi\)）のエネルギーが時間的に変化せずピタリと安定している場合、方程式は次のような極めて美しいシンプルな形（時間に依存しないシュレディンガー方程式）に落ち着きます。

\[\hat{H} \psi = E \psi\]

あるオペレーター（\(\hat{H}\)）を、特定の綺麗な形をした波\(\psi\)に作用させたとき、波の形は全く崩れず、ただ前に「単なる数字\(E\)」がポロッと飛び出してくる。

この特別な波の形\(\psi\)を\textbf{「固有関数（固有状態）」と呼び、飛び出してきた数字\(E\)を「固有値」}と呼びます。

この数学的な関係は、物理学的にとてつもない意味を持っています。この飛び出してきた「固有値」こそが、私たちが実験室のメーターで読み取る「現実の観測データ（観測可能量：オブザーバブル）」そのものなのです。

これは、楽器がどんな音でも連続的に出すのではなく、その形に合った共鳴音を特に強く鳴らすことに少し似ています。原子の中の電子波も、原子核の周りで許される「共鳴の形」だけが安定に存在でき、そのときのエネルギーだけが固有値として現れます。

この数式に水素原子の条件（原子核からの引力\(V\)）を入れて解くと、第8章でボーアがデータから無理やり「飛び飛びだ」と仮定していた電子の軌道やエネルギー準位が、「方程式が特定の固有値しか持たない（波が空間に綺麗に収まる定在波の条件）」という極めて自然な数学的帰結として自動的に導き出されたのです。

\section{行列力学の誕生とハイゼンベルクの不確定性原理}

シュレディンガーが波の数式を完成させたのと全く同時期の1925年、ドイツの天才ヴェルナー・ハイゼンベルクも、全く別のアプローチから量子力学を完成させていました。それが\textbf{「行列力学」}です。

\subsection{見えない「軌道」を捨て、観測できる「データ」だけを使う}

ハイゼンベルクの出発点は、ボーアの原子模型に対する強い疑念でした。ボーアの模型は「電子が特定の軌道を回っている」と仮定していましたが、誰もその軌道を実際に見たことはありません。ハイゼンベルクは、「観測不可能な『軌道』のようなフィクション（作り話）に頼るのをやめ、実験で確実に観測できるデータ（原子が放つ光の振動数や明るさ）だけで物理学を再構築すべきだ」という、極めて現代的でデータ駆動的なアプローチをとりました。

彼は、電子の位置や運動量を単なる「数字」ではなく、無数の観測データを縦横に規則正しく並べた\textbf{「行列（マトリックス）」}という数学の表として表現しました。

\subsection{行列の「固有値」が現実を決める}

ここで奇跡が起きます。ハイゼンベルクが構築した巨大な「エネルギーの行列」に対して、線形代数学における\textbf{「固有値」を求める計算を行いました（行列の対角化）。すると、その計算結果として出てきた無数の固有値が、現実の水素原子が放つ光のエネルギーの飛び飛びのデータと完全に一致}したのです。

シュレディンガーの「微分オペレーターの固有値」も、ハイゼンベルクの「行列の固有値」も、どちらも\textbf{「観測される現実の値は、数学的な構造の中に潜む『固有値』として現れる」}という全く同じ真理を突いていました。

（※後に、ハイゼンベルクの「行列力学」とシュレディンガーの「波動力学」は、見た目の数式は全く違うものの、数学的には完全に等価なものであることが証明されました。）

\subsection{行列の「掛け算の順番」が不確定性を生み出す}

さらにハイゼンベルクは、物理量を「行列」や「オペレーター」にしたことによる、ある決定的な数学的性質に気づきました。それが\textbf{「行列の掛け算は、かける順番を逆にすると答えが変わってしまう（非可換性）」}ということです。

ただの数字なら、\(3 \times 5 = 5 \times 3\)です。しかし行列（やシュレディンガーのオペレーター）は違います。「靴下を履いてから靴を履く」のと「靴を履いてから靴下を履く」のでは結果が全く異なるように、量子力学の世界において\textbf{「位置を表す行列（\(\hat{x}\)）」と「運動量を表す行列（\(\hat{p}\)）」を掛け合わせる順番を入れ替えると、結果に必ずズレ（\(i\hbar\)の差）が生じてしまう}ことが数学的に証明されたのです。

\[[\hat{x},\hat{p}] = \hat{x}\hat{p} - \hat{p}\hat{x} = i\hbar\]

この「順番によるズレ」こそが、ミクロの世界における最も不気味な性質\textbf{「不確定性原理」}の真の正体です。

「電子の『位置』と『運動量（スピード）』は、同時に正確に測ることは絶対にできない」。位置を正確に知ろうとするとスピードが全く分からなくなり、スピードを正確に測ろうとすると位置がぼやけてしまいます。これは人間の観測技術が未熟だからではなく、宇宙の法則（行列の非可換性）が「両方を同時に確定させることを禁じている」からなのです。数式では、位置の不確かさ\(\Delta x\)と運動量の不確かさ\(\Delta p\)の間に、次の下限が現れます。

\[
\Delta x\,\Delta p \geq \frac{\hbar}{2}
\]

ここに至って、ニュートンやラプラスが夢見た「現在の状態が完璧に分かれば、未来は計算で完璧に予測できる」という決定論的宇宙観は完全に崩壊しました。

\section{ボルンの確率解釈：決定論の終焉と「波の収縮」}

ハイゼンベルクの行列力学とシュレディンガーの波動方程式は、見た目の数学は全く違うものの、同じエネルギーの「固有値」を導き出すことが数学的に証明されました。しかし、ここで物理学者たちを真っ二つに分断する、科学史上の最大の大論争が巻き起こります。

「シュレディンガー方程式が計算している\textbf{『波（波動関数\(\psi\)）』とは、現実の宇宙において一体何を表しているのか？}」という問題です。

\subsection{シュレディンガーの誤算：「実在の波」ではない}

方程式を作ったシュレディンガー自身は、電子という物質そのものが、空間にモヤモヤとゼリーのように広がった「実在の波（電荷の雲）」として存在していると考えていました。

しかし、この考え方には致命的な無理がありました。もし電子が本当に空間に広がった物理的な波なら、それを半分に割ったり、少しずつ削り取ったりできるはずです。しかし実際の実験では、電子は常に「1個、2個」と数えられる「完全な硬い粒」としてしか観測されません。観測器に当たるのは、常に空間の「ただ一点」なのです。

\subsection{なぜ「絶対値の2乗」なのか？ 複素数から実数への変換}

この矛盾を解決し、現代の量子力学の正解となる解釈を打ち立てたのが、ドイツの物理学者マックス・ボルンです。彼は1926年、次のような驚くべき宣言を行いました。

「波動関数\(\psi\)それ自体は、目に見える実在の波ではない。その振幅の絶対値の2乗（\(|\psi|^2\)）が、\textbf{『そこに電子が粒として発見される確率』}を表しているに過ぎない」

なぜわざわざ「絶対値の2乗（ノルムの2乗）」を計算するのでしょうか？

10.2節で見たように、波動関数\(\psi\)は「虚数\(i\)」を含んだ複素数の波です。現実世界で「明日の降水確率は\(30 + 10i\)％です」と言われても全く意味が分からないように、確率は必ず「0から1の間の実数」でなければなりません。

数学において、複素数の絶対値の2乗（自分自身とその複素共役を掛け合わせたもの）を計算すると、虚数部分が見事に消え去り、必ず「0以上の実数」になります。ボルンは、目に見えない複素数の波から、私たちが観測できる現実の「確率」という実数を紡ぎ出すための数学的ルールとして、この\(|\psi|^2\)を導入したのです。

\subsection{コペンハーゲン解釈と「波束の収縮」}

このボルンの「確率解釈」を基盤として、ニールス・ボーアやハイゼンベルクらが中心となってまとめ上げた量子力学の公式見解を\textbf{「コペンハーゲン解釈」}と呼びます。これは、私たちの常識を根本から破壊する恐ろしい世界観でした。

電子は、人間が「観測」する前は、どこにいるか決まっておらず「ここにも、あそこにもいる」という複素数の確率の波として空間全体に重なり合って広がっています（重ね合わせの状態）。

しかし、測定器と相互作用して「観測された」瞬間に、この広がっていた波は一つの結果へ絞り込まれたように見え（波束の収縮）、サイコロが振られたように確率に従って空間の「ただ一点」に粒として現れるのです。ここでいう観測とは、人間の意識が見つめることではなく、電子の状態が測定装置や周囲の環境と切り離せないほど強く関わることだと考えると理解しやすくなります。

月は私たちが見ていなくてもそこに存在します。しかし量子力学の世界では、少なくとも電子の位置のようなミクロな量について、\textbf{「観測するまでは一つの値に確定しているとは言えず、観測によって一つの結果が現れる」}ことになります。そして、次にサイコロの目がどう出るか（どこで電子が観測されるか）を事前に完全には予測できません。予測できるのは、\(|\psi|^2\)が与える確率分布までです。

\subsection{思考実験の極致：「シュレディンガーの猫」}

この「観測するまで現実が重なり合っている」という奇妙な解釈に、シュレディンガー本人は猛烈に反発しました。「自分が作った美しい波の数式が、サイコロ遊びのような確率に堕ちてしまった」と嘆いた彼は、コペンハーゲン解釈の不条理さを浮き彫りにするため、1935年に世にも奇妙な思考実験を提唱しました。それが\textbf{「シュレディンガーの猫」}です。

密閉された鋼鉄の箱の中に、猫を1匹入れます。

箱の中には「1時間以内に50\%の確率で崩壊して放射線を出す原子」と、「放射線を検知すると毒ガスが出る装置」をセットします。

1時間後、原子が崩壊していれば毒ガスが出て猫は「死んで」おり、崩壊していなければ猫は「生きて」います。

コペンハーゲン解釈に忠実に従えば、箱を閉めている間、ミクロな原子は「崩壊している状態」と「崩壊していない状態」が50\%ずつ重なり合っています。となれば、それと連動しているマクロな猫もまた、\textbf{「生きている状態」と「死んでいる状態」が重なり合った、日常感覚では受け入れがたい状態（重ね合わせ）}で箱の中に存在していることになります。

そして、人間が箱のフタを開けて「観測した」瞬間に波が収縮し、生きているか死んでいるかの「どちらか」に現実が確定するのです。

シュレディンガーは、「ミクロの確率的な曖昧さが、マクロな猫の生死にまで拡大解釈されるなんて馬鹿げている！」と批判しましたが、皮肉なことに、現在ではこの思考実験こそが「量子力学の重ね合わせ」を説明する最も有名な例え話として定着しています。

\subsection{アインシュタインの最後の抵抗：「量子もつれ」とEPRパラドックス}

このコペンハーゲン解釈の曖昧で主観的な世界観に対し、最後まで最も強力に抵抗したのがアインシュタインでした。

彼は\textbf{「神はサイコロを振らない」「私が見ていなくても月はそこにあるはずだ」}と言って、量子力学にはまだ私たちの知らない「隠れた変数（決定論的なメカニズム）」があるはずだと主張しました。

1935年、アインシュタインはポドルスキー、ローゼンと共に、量子力学の致命的な矛盾を突く\textbf{「EPRパラドックス」という論文を発表しました。 彼らが目をつけたのは、ペアで生まれた2つの粒子が持つ「量子もつれ（エンタングルメント）」}という現象です。2つの粒子AとBがもつれ状態にあるとき、片方（A）のスピン（自転）が「上向き」と観測されると、もう片方（B）は必ず「下向き」になるという強い相関関係を持っています。

アインシュタインは思考実験を行いました。もつれ合った粒子AとBを、宇宙の端と端（何光年も離れた場所）に引き離します。

もし量子力学が言うように「観測するまで状態が決まっていない」のなら、Aを観測して「上向き」と確定した瞬間に、何光年も離れたBに対して「お前は下向きになれ！」という情報が「一瞬で（光の速さを超えて）」伝わらなければなりません。

アインシュタインはこれを\textbf{「不気味な遠隔作用（Spooky action at a distance）」}と呼び、「相対性理論によれば光より速く伝わる情報はないのだから、こんな現象はあり得ない。粒子AとBは、生まれたときから『上と下になる』とあらかじめ決定されていた（隠れた変数があった）と考えるべきだ」と結論づけました。

しかし、その後の物理学の歴史は、アインシュタインにとって残酷なものでした。

1964年にジョン・ベルが考案した「ベルの不等式」を用いた1980年代以降の精密な実験（アラン・アスペらの実験、2022年ノーベル物理学賞）によって、\textbf{「局所的な隠れた変数だけでは、量子もつれの相関を説明できない」}ことが示されました。つまり、粒子AとBが単に「出発時に答えを書いた封筒を持っていた」だけ、という古典的な説明では足りないのです。

ただし、量子もつれを使って光より速くメッセージを送れるわけではありません。離れた粒子の測定結果には強い相関がありますが、個々の結果そのものは確率的で、こちらの意志で好きな情報を相手へ送ることはできないからです。アインシュタインの願いは退けられ、自然が古典的な意味での局所実在論には従わないことが明らかになりました。そして、この「量子もつれ」こそが、量子コンピュータや量子通信など、現代の量子情報科学の中心的な資源になっています。

\section{量子力学の絶望的な壁：「次元の呪い」とスパコンの限界}

こうして1920年代に完成した量子力学は、原理的には、この宇宙にある多くの物質現象（原子、分子、化学反応、薬の効き目、新素材の性質など）を計算し、予測できるはずのものでした。イギリスの天才物理学者ポール・ディラックは、「物理学と化学の大部分を数学的に計算するための基本法則は、完全に解明された」と高らかに勝利宣言しました。

しかし、ディラックは同時にこうも付け加えました。「ただ一つの問題は、その方程式が複雑すぎて、人間の手には負えないことだ」。

実は、量子力学の計算には\textbf{「次元の呪い（Curse of Dimensionality）」}と呼ばれる、絶望的な数学の壁が立ちはだかっていました。

\subsection{波動関数における「指数関数的」な計算量の爆発}

水素原子のように、電子が1個しかない場合は「紙と鉛筆」でも計算できます。しかし、水分子（H\(_2\)O）やタンパク質、あるいは半導体のように「複数の電子」が互いに反発し合いながら存在するシステム（多体問題）になると、状況は一変します。

シュレディンガーの波動関数で計算しようとした場合を考えてみましょう。電子1個の位置を表すために、空間を100個のマス目（メッシュ）に分けたとします。電子が2個になると、お互いの位置の組み合わせをすべて計算しなければならないため、\(100 \times 100 = 10,000\)個のパラメータが必要になります。電子が\(N\)個になれば、必要な変数の数は\(100^N\)という指数関数的（爆発的）な勢いで増大していきます。

\subsection{行列力学とヒルベルト空間の「絶望的な広がり」}

これは、ハイゼンベルクの「行列力学」を使っても全く同じ問題に行き着きます。

量子力学において、物質が取りうる「すべての状態（重ね合わせのパターン）」を数学的に表現するための仮想的な多次元空間を\textbf{「ヒルベルト空間（Hilbert Space）」}と呼びます。

私たちが生きる現実の空間は縦・横・高さの3次元ですが、ヒルベルト空間は異なります。そこでは「状態のパターンの数」そのものが、空間の新しい「軸（次元）」となります。そのため、粒子が一つ増えたり、状態の選択肢が増えたりするだけで、ヒルベルト空間の次元数は掛け算式に爆発的に増えていくのです。

例えば、電子が自転のような性質（スピン）を持っており、それぞれ「上向き」と「下向き」の2つの状態をとれるとします。電子が10個あるだけで、全体の組み合わせの数、すなわちヒルベルト空間の次元数は\(2^{10} = 1024\)次元になります。

もし電子が300個集まった物質（例えば、ごく小さな分子）を計算しようとすると、ヒルベルト空間の次元は\(2^{300}\)次元となります。驚くべきことに、この数字は、私たちが観測できる宇宙全体に存在するすべての原子の数（約\(10^{80}\)個）よりもはるかに大きいのです。ハイゼンベルクのルールに従ってこの物質のエネルギーを計算しようとすれば、宇宙の原子の数よりも多くの縦横のマス目を持つ「超巨大な行列」の固有値を求めなければならず、それは原理的に不可能でした。

\subsection{スーパーコンピュータの完全なる敗北}

この「次元の呪い」の前では、現代の最新鋭のテクノロジーも全く無力です。「富岳」のような世界最速のスーパーコンピュータを何万台つなぎ合わせ、宇宙の誕生から終わりまでの時間を使って計算させたとしても、たった数十個の電子の厳密な動きすら解き明かすことはできません。

コンピュータの性能が1年半で2倍になるという「ムーアの法則」をもってしても、計算対象の分子の電子が「たった1個」増えるだけで、その進化分は一瞬で相殺されてしまいます。ディラックの勝利宣言から約100年、量子力学の方程式を解いて「新しい薬」や「夢の素材」を自由に設計することは、計算量の壁によって阻まれ続けてきました。

\section{AIが解く量子力学：ニューラルネットワークと波動関数の融合}

この絶望的な「次元の呪い」を打ち破るために、物理学者たちは様々な「近似手法（大雑把な計算で妥協する方法）」を開発してきましたが、どうしても精度に限界がありました。しかし2010年代後半、この現代物理学の最大の壁に、ついに\textbf{AI（ディープラーニング）}が挑み始め、数々のブレイクスルーを起こしています。

\subsection{ブレイクスルー①：波動関数をAIに丸暗記させる（NNQS）}

2017年、Giuseppe Carleo と Matthias Troyer という二人の研究者が、\textbf{「Neural Network Quantum States (NNQS)」}という画期的な手法を発表しました。彼らは、シュレディンガー方程式の答えである複雑怪奇な「波動関数」そのものを、人工知能の「ニューラルネットワーク」で表現させてしまおうと考えたのです。

ニューラルネットワークは、人間の顔や言葉などの「極めて複雑で高次元なパターン」を、比較的少ないパラメータで効率よく圧縮・表現する天才です（万能近似定理）。

研究者たちは、AIに「システムの総エネルギー（ハミルトニアンの期待値）をできるだけ低くするように、自分のネットワークの重みを調整しなさい」という物理学のルール（変分原理）だけを与えました。するとAIは、膨大なパズルを解くように自己学習を進め、指数関数的に大きな波動関数をそのまま保存することなく、重要な構造だけを圧縮して表すことに成功したのです。

特に、無数の電子が複雑に絡み合う「量子もつれ（エンタングルメント）」のパターンを、AIの隠れ層のネットワーク構造が見事に模倣・圧縮できることが判明し、物理学界に衝撃を与えました。

\subsection{ブレイクスルー②：パウリの排他原理を学習する「FermiNet」}

しかし、複数の電子を扱う量子化学の計算には、さらにもう一つ、AIにとって極めて厄介な物理法則の壁がありました。それが\textbf{「パウリの排他原理」}です。

電子は「フェルミ粒子」と呼ばれるグループに属しており、「2つの電子は、全く同じ状態（同じ場所、同じスピン）をとることができない」という強い反発のルールを持っています。これを波動関数の数学で表すと、「電子Aと電子Bの位置を入れ替えると、波のプラスマイナス（符号）が完全に反転しなければならない（反対称性）」という非常に厳しい制約になります。

通常のニューラルネットワークに電子の位置を入力しても、この「入れ替えると符号が反転する」という奇妙な性質をうまく表現できませんでした。

この難問を解決したのが、2020年にGoogle DeepMindが発表した\textbf{「FermiNet（フェルミネット）」}です。彼らは、AIのネットワークの出力層に「行列式（スレイター行列式）」という数学的なフィルターを強制的に組み込みました。行列式は、行や列を入れ替えると必ず符号が反転するという性質を持っています。

これにより、FermiNetは「パウリの排他原理を本能として絶対に守るAI」となり、従来の近似計算では到達できなかった極めて高い精度で、複数の電子が複雑に絡み合う分子のエネルギーを計算することに成功しました。

\subsection{ブレイクスルー③：未知の数式をAIで補完する「DM21」}

現在、世界中の材料科学や創薬の現場で最も使われている量子力学の計算手法に\textbf{「密度汎関数理論（DFT）」があります。これは、次元の呪いを避けるために、多電子の複雑な波を「空間の電子の密度（3次元）」という単純な雲のようなものに置き換えて計算する素晴らしい手法です。 しかしDFTには、「電子同士がどう避け合うか」という効果を表す「交換相関汎関数」という部分の正確な数式が、人類には未だに分かっていない}という致命的な弱点がありました。物理学者たちは何十年もかけて、この未知の数式を人間の直感で推測（近似）してきましたが、限界に達していました。

ここでもDeepMindが革命を起こします。彼らは2021年に\textbf{「DM21」}というAIモデルを発表しました。これは、高精度な実験データや厳密な計算データを用いて、「未知の交換相関汎関数の数式（パズルのピース）」そのものを、ニューラルネットワークに学習・推論させるというアプローチです。

DM21は、いくつかの重要なベンチマークで従来の汎関数を上回る精度を示しました。AIはもはや計算を速くするだけでなく、「人類が解き明かせなかった量子力学の方程式の欠落部分」を、データから補完する存在になりつつあるのです。

\subsection{ブレイクスルー④：超高速な「仮想の実験室」を作るニューラルネットワークポテンシャル}

さらに、量子力学の計算結果をAIに学習させ、「原子がどう動くかのシミュレーション（分子動力学法）」をスパコンの数百万倍のスピードで実行する技術が急速に普及しています。

本来、水が凍る過程や、タンパク質が薬と結合する様子をシミュレーションするには、1コマ（1フェムト秒）進むごとに、すべての電子と原子核についてシュレディンガー方程式を解き直さなければなりません。これはスパコンを使っても数ヶ月〜数年かかる絶望的な計算です。

そこで研究者たちは、あらかじめ様々な原子の配置における「量子力学的な力の働き方」を厳密に計算し、その膨大なデータをAIに学習させました。するとAIは、「この配置なら、原子はこっちにこれくらいの力で引っ張られるはずだ」という量子力学の答えを、方程式を解かずに瞬時に『推論』できるようになります（ニューラルネットワークポテンシャル：NNP）。

このAIの直感的な推論を使うことで、毎ステップで重い量子計算を解き直す負担を大きく減らし、多数の原子が織りなす複雑な化学反応や物質の相転移を、より現実的な時間でシミュレーションできる「仮想実験室」が現実のものになりつつあるのです。

\subsection{物質科学と創薬のルネサンス}

これらのAIと量子力学の融合技術は、いまや現実世界に多大な恩恵をもたらし始めています。

従来のスパコンでは計算不可能だった複雑なタンパク質の折り畳み構造（アルツハイマー病などの原因）の解明や、送電ロスをゼロにする常温超伝導物質の探索、そして効率的な全固体バッテリーの材料設計など。AIは、ディラックが残した「解けない方程式」の次元の壁をすり抜け、コンピュータ上で未知の物質を自由自在にデザインする「魔法の杖」として、現代科学にルネサンスをもたらそうとしています。

\section*{【第10章 章末ディスカッション課題】}

量子力学は「未来は確率でしか決まらない」ことを示しました。もし人間の脳の思考や意識も、脳内の電子の量子力学的な動きによって生まれているとすれば、私たちには「自由意志」があると言えるでしょうか？ それとも、すべてはサイコロの確率の産物なのでしょうか。


\section*{参考文献（学術誌）}
\begin{enumerate}[leftmargin=2.4em]
\item Heisenberg, W. ``Uber den anschaulichen Inhalt der quantentheoretischen Kinematik und Mechanik.'' \textit{Zeitschrift fur Physik}, 43, 172--198, 1927. DOI: \url{https://doi.org/10.1007/BF01397280}.
\item Bell, J. S. ``On the Einstein Podolsky Rosen paradox.'' \textit{Physics Physique Fizika}, 1(3), 195--200, 1964. DOI: \url{https://doi.org/10.1103/PhysicsPhysiqueFizika.1.195}.
\item Aspect, A., Dalibard, J., and Roger, G. ``Experimental test of Bell's inequalities using time-varying analyzers.'' \textit{Physical Review Letters}, 49(25), 1804--1807, 1982. DOI: \url{https://doi.org/10.1103/PhysRevLett.49.1804}.
\item Carleo, G. and Troyer, M. ``Solving the quantum many-body problem with artificial neural networks.'' \textit{Science}, 355(6325), 602--606, 2017. DOI: \url{https://doi.org/10.1126/science.aag2302}.
\item Pfau, D., Spencer, J. S., Matthews, A. G. D. G., and Foulkes, W. M. C. ``Ab initio solution of the many-electron Schrodinger equation with deep neural networks.'' \textit{Physical Review Research}, 2, 033429, 2020. DOI: \url{https://doi.org/10.1103/PhysRevResearch.2.033429}.
\end{enumerate}


